\chapter{Manual de usuario}
\label{anexo:manual_usuario}

Este anexo recoge la guía de uso de OmniLitter desde el punto de vista de las
personas que van a utilizar el sistema. Su objetivo es explicar, de forma sencilla,
cómo moverse por el portal web y por la aplicación móvil, qué pasos seguir en los
flujos habituales y qué significa cada apartado de la plataforma.

El manual no describe detalles internos de implementación. Se centra en las
acciones que realizará el usuario final: entrar en la plataforma, seleccionar su
grupo de trabajo, consultar resultados, gestionar campañas y registrar datos en
campo.

\newcommand{\capturapendiente}[2]{
\begin{figure}[H]
  \centering
  \fbox{
    \begin{minipage}[c][0.22\textheight][c]{0.82\textwidth}
      \centering
      \textbf{Captura pendiente}\\[0.4em]
      #1
    \end{minipage}
  }
  \caption{#2}
\end{figure}
}

\newcommand{\capturamanual}[3]{
\begin{figure}[H]
  \centering
  \includegraphics[width=#3\textwidth]{#1}
  \caption{#2}
\end{figure}
}

\section{Conceptos básicos}

OmniLitter organiza el trabajo en torno a grupos, campañas, sesiones de muestreo
y observaciones.

\begin{description}
  \item[Grupo.] Equipo de trabajo al que pertenece el usuario. Cada grupo tiene sus
  propios miembros, campañas y datos. Antes de trabajar, el usuario debe comprobar
  cuál es su grupo activo.
  \item[Campaña.] Actividad de muestreo planificada para una zona, periodo o
  propósito concreto. Por ejemplo, una campaña de primavera en una playa o en una
  ribera.
  \item[Sesión.] Jornada o recorrido concreto dentro de una campaña. La sesión
  indica el entorno, el método de muestreo y el recorrido seguido.
  \item[Observación.] Registro de un residuo o conjunto de residuos encontrado en
  campo. Puede incluir ubicación, tipo de residuo, cantidad, notas y fotografías.
\end{description}

La plataforma contempla tres perfiles de uso:

\begin{itemize}[leftmargin=1.5em]
  \item \textbf{Investigador de campo}: registra observaciones desde la aplicación
  móvil y consulta los datos a los que tiene acceso.
  \item \textbf{Administrador de grupo}: además de lo anterior, puede gestionar
  miembros del grupo, crear campañas y mantener organizada la información del equipo.
  \item \textbf{Superadministrador}: puede consultar y administrar la plataforma de
  forma global, incluyendo todos los grupos activos.
\end{itemize}

Las opciones visibles pueden variar según el perfil. Si una acción no aparece en
pantalla, normalmente significa que el usuario no tiene permisos para realizarla en
ese grupo.

\section{Acceso a OmniLitter}

\subsection{Entrada desde el portal web}

Al abrir el portal web, el usuario llega a una página inicial desde la que puede
iniciar sesión o crear una cuenta. Para entrar, debe pulsar \textit{Iniciar sesión},
introducir su correo electrónico y contraseña, y confirmar el acceso.

\capturamanual{figures/pagina-inicial-web.png}{Pantalla de bienvenida del portal web.}{0.95}

Si el usuario todavía no tiene cuenta, puede registrarse desde \textit{Crear cuenta}.
Si ha olvidado su contraseña, debe usar \textit{¿Olvidaste tu contraseña?} y seguir
el proceso de recuperación indicado por la plataforma.

\capturamanual{figures/inicio-sesion-web.png}{Acceso de usuario al portal web.}{0.82}

\subsection{Entrada desde la aplicación móvil}

La aplicación móvil utiliza las mismas credenciales que el portal web. En la pantalla
de inicio de sesión se introduce el correo electrónico y la contraseña. Si el
dispositivo no tiene conexión, la app permite entrar solo si esa cuenta ya se había
utilizado antes con conexión en el mismo dispositivo.

\capturamanual{figures/inicio-sesion-movil.png}{Acceso de usuario en la aplicación móvil.}{0.42}

Cuando se trabaja en campo, es recomendable iniciar sesión con conexión antes de
salir. Así la aplicación puede preparar los datos necesarios para trabajar aunque
después se pierda la cobertura.

\section{Uso del portal web}

El portal web está pensado para consultar, revisar, analizar y administrar la
información recogida en campo. Tras iniciar sesión, se muestra un menú lateral con
los apartados principales y, en la parte superior del menú, el usuario activo y el
grupo de trabajo seleccionado.

\capturamanual{figures/estructura-web.png}{Estructura principal del portal web.}{0.95}

\subsection{Seleccionar el grupo activo}

El selector de grupo determina qué datos se ven en la mayoría de pantallas. Si el
usuario pertenece a varios grupos, debe elegir el grupo con el que quiere trabajar
antes de consultar campañas, observaciones, mapas o métricas.

En el caso del superadministrador, dejar el selector sin un grupo concreto permite
ver una vista agregada de todos los grupos activos. Para el resto de usuarios, lo
normal es trabajar siempre con un grupo activo.

\subsection{Dashboard}

El \textit{Dashboard} ofrece una vista rápida del estado del grupo: número de
observaciones, residuos registrados, fotografías, sesiones, campañas, evolución
temporal y rankings principales.

Para consultar un periodo concreto, el usuario puede usar los campos \textit{Desde}
y \textit{Hasta}. Si quiere volver a la vista completa, debe pulsar \textit{Limpiar}.
Cuando hay datos disponibles, el panel permite exportar la información en CSV o JSON.

Este apartado resulta útil para responder preguntas rápidas, por ejemplo cuántas
observaciones se han registrado en una campaña, qué materiales aparecen con más
frecuencia o cómo evoluciona la actividad del grupo.

\subsection{Campañas}

La pantalla \textit{Campañas} muestra las campañas del grupo activo. El usuario puede
buscar por nombre, filtrar por estado y consultar las observaciones asociadas a cada
campaña.

\capturamanual{figures/campañas-web.png}{Pantalla de campañas del portal web.}{0.95}

Los administradores pueden crear una nueva campaña desde \textit{Nueva campaña}. El
formulario solicita el nombre, una descripción opcional, el estado inicial y las
fechas de inicio y fin. Una campaña nueva puede quedar como borrador o activarse
directamente.

\capturamanual{figures/editar-campaña.png}{Alta y edición de campañas.}{0.86}

En el listado, los administradores pueden editar o eliminar campañas. Los
investigadores de campo tienen una vista más orientada a la consulta: pueden ver sus
campañas y otras campañas del grupo, pero no modificar la planificación si no tienen
permisos de administración.

\subsection{Observaciones}

La pantalla \textit{Observaciones} permite consultar los registros creados durante el
trabajo de campo. Cada fila muestra el tipo de residuo, la cantidad, la sesión y
campaña asociadas, el entorno, el autor, la fecha y si tiene fotografías.

\capturamanual{figures/observaciones-web.png}{Consulta de observaciones en el portal web.}{0.95}

El usuario puede filtrar por campaña, sesión, entorno, autor, presencia de fotos y
rango de fechas. También puede usar el buscador para localizar una observación por
tipo de residuo, nombre de campaña, sesión, autor o nota.

Al pulsar una observación se abre su detalle. En esta vista se consultan las
coordenadas, el observador, los elementos registrados, las notas y las fotografías
asociadas. Las fotografías pueden ampliarse para revisarlas con más detalle.

\capturamanual{figures/detalle-observacion-web.png}{Detalle de observación y evidencias fotográficas.}{0.9}

\subsection{Mapa}

El apartado \textit{Mapa} muestra las observaciones geolocalizadas. Dispone de tres
modos:

\begin{itemize}[leftmargin=1.5em]
  \item \textbf{Marcadores}: cada punto representa una observación. Al pulsarlo se
  muestran sus datos básicos y, si existe, una fotografía.
  \item \textbf{Densidad}: muestra las zonas con mayor concentración de residuos.
  \item \textbf{Rutas}: representa los recorridos de las sesiones de muestreo, las
  áreas cubiertas y las observaciones tomadas durante el recorrido.
\end{itemize}

\capturamanual{figures/mapa-web.png}{Mapa de observaciones y rutas de muestreo.}{0.95}

Los filtros permiten acotar por fechas, campaña, entorno y categoría de residuo. Si
se necesita trabajar con los datos en otra herramienta, el mapa permite exportar el
resultado en formato GeoJSON.

\subsection{Analítica}

La pantalla \textit{Analítica} ofrece una vista más detallada que el dashboard. Está
pensada para explorar los datos antes de preparar informes o exportaciones.

\capturamanual{figures/analititca-web.png}{Analítica avanzada del portal web.}{0.95}

El usuario puede filtrar por campaña, sesión, entorno, autor y fechas. A partir de
esos filtros, el portal muestra indicadores resumidos, un ranking de materiales, un
desglose por categorías y una tabla con el detalle de las observaciones.

Al pulsar una barra del ranking o una categoría, la tabla inferior se centra en ese
subconjunto. Esta interacción ayuda a revisar de forma rápida de dónde vienen los
datos que aparecen en las gráficas.

Desde \textit{Exportar}, el usuario puede descargar observaciones, detalle de
residuos, resúmenes y ficheros preparados para análisis científico. Las exportaciones
científicas incluyen un diccionario de columnas para facilitar la interpretación de
los datos fuera de OmniLitter.

\subsection{Grupos}

El apartado \textit{Grupos} permite consultar y administrar los equipos de trabajo.
El listado separa los grupos activos de los archivados e incluye buscador, filtro de
estado y opciones de ordenación.

\capturamanual{figures/grupos-web.png}{Gestión de grupos de trabajo.}{0.95}

Al abrir un grupo se puede revisar su descripción, miembros y campañas asociadas. Los
administradores pueden editar el grupo, añadir miembros, enviar invitaciones,
archivar el grupo o restaurarlo si estaba archivado.

\capturamanual{figures/detalles-grupos-web.png}{Detalle y miembros de un grupo.}{0.9}

Un usuario también puede abandonar un grupo si ya no participa en él. Si lo abandona
por error, el portal ofrece durante unos minutos la opción de volver a unirse.

\subsection{Usuarios}

La pantalla \textit{Usuarios} muestra las cuentas de la plataforma o del grupo activo,
según el perfil del usuario. El superadministrador puede crear usuarios, activar o
desactivar cuentas, cambiar roles y gestionar la pertenencia a grupos.

\capturamanual{figures/usuarios-web.png}{Gestión de usuarios en el portal web.}{0.95}

Los administradores de grupo tienen una vista centrada en los miembros de su grupo.
Desde ella pueden invitar usuarios, cambiar el rol dentro del grupo y expulsar a
miembros cuando sea necesario. Los investigadores de campo pueden consultar la
información disponible, pero no realizar acciones de gestión.

\subsection{Configuración}

En \textit{Configuración} el usuario puede actualizar su perfil, cambiar el idioma de
la interfaz, modificar la contraseña y revisar invitaciones pendientes a grupos.

\capturamanual{figures/configuracion-web.png}{Configuración del usuario en el portal web.}{0.95}

También se puede elegir el formato de exportación preferido. Los administradores de
grupo pueden descargar una copia administrativa de su grupo, y el superadministrador
puede descargar una copia global de los grupos activos. El superadministrador cuenta,
además, con un apartado para gestionar el catálogo de entornos y tipologías usado en
las campañas.

\section{Uso de la aplicación móvil}

La aplicación móvil está pensada para el trabajo en campo. Permite preparar campañas,
abrir sesiones de muestreo, registrar observaciones con GPS, añadir fotografías y
sincronizar los datos cuando haya conexión.

La navegación principal está en la parte inferior de la pantalla y contiene cinco
apartados: \textit{Inicio}, \textit{Campañas}, \textit{Mapa}, \textit{Sincronización}
y \textit{Perfil}.

\capturamanual{figures/inicio-navegacion-movil.png}{Navegación principal de la aplicación móvil.}{0.42}

\subsection{Inicio}

La pantalla de inicio muestra el grupo activo, el estado de conexión, las sesiones
recientes y un acceso directo a las campañas. También indica si existen datos
pendientes de sincronizar.

\capturamanual{figures/inicio-navegacion-movil.png}{Pantalla de inicio de la aplicación móvil.}{0.42}

Si el dispositivo está sin conexión, aparece un aviso indicando que los datos se
guardarán localmente. El usuario puede continuar trabajando y sincronizar más tarde.

\subsection{Perfil y grupos}

Desde \textit{Perfil} se revisan los datos personales, el rol del usuario y el grupo
activo. El usuario puede seleccionar otro grupo, trabajar sin grupo activo, crear un
grupo nuevo o aceptar invitaciones pendientes.

\capturamanual{figures/perfil-movil.png}{Perfil de usuario en la aplicación móvil.}{0.42}
\capturamanual{figures/grupos-movil.png}{Selección y gestión de grupos en la aplicación móvil.}{0.42}

Antes de registrar datos, es importante comprobar que el grupo activo es el correcto.
Las campañas, sesiones y observaciones quedarán asociadas a ese contexto de trabajo.

Los administradores de grupo pueden invitar usuarios desde el perfil. Para enviar o
responder invitaciones se necesita conexión, porque la invitación debe validarse en
la plataforma.

\subsection{Campañas en la app móvil}

El apartado \textit{Campañas} reúne el flujo completo de trabajo en campo. Desde esta
pantalla el usuario puede seleccionar una campaña, crear una sesión y registrar
observaciones.

\capturamanual{figures/campañas-movil.png}{Campañas en la aplicación móvil.}{0.42}

Si el usuario tiene permisos de administración, puede crear campañas desde la app.
El formulario solicita el nombre, la descripción, el entorno, la fecha de inicio y
el estado. Si no hay conexión, la campaña puede quedar guardada localmente hasta que
se sincronice.

\capturamanual{figures/crear-campaña.png}{Creación de campaña desde la aplicación móvil.}{0.42}
\capturamanual{figures/crear-campaña2.png}{Detalle del formulario de campaña en la aplicación móvil.}{0.42}

\subsection{Crear y finalizar una sesión de muestreo}

Dentro de una campaña, el usuario puede abrir una nueva sesión de muestreo. Para
crear una sesión debe indicar un nombre, el entorno, la tipología si aplica y el
método de muestreo.

La aplicación contempla tres métodos:

\begin{itemize}[leftmargin=1.5em]
  \item \textbf{Libre}: recorrido abierto sin una forma predeterminada.
  \item \textbf{Transecto}: recorrido lineal con un ancho de banda definido.
  \item \textbf{Cuadrícula}: muestreo por celdas regulares.
\end{itemize}

\capturamanual{figures/seleccion-entorno-movil.png}{Selección del entorno de muestreo en la aplicación móvil.}{0.42}
\capturamanual{figures/seleccion-tipologia-movil.png}{Selección de tipología de muestreo en la aplicación móvil.}{0.42}
\capturamanual{figures/crear-sesion-movil.png}{Creación de una sesión de muestreo.}{0.42}

Al crear la sesión, la app solicita permiso de ubicación para guardar el punto GPS de
inicio. Durante la sesión, el usuario puede consultar el mapa del recorrido y añadir
observaciones. Al terminar el trabajo de campo debe pulsar \textit{Finalizar sesión}.
La sesión quedará marcada como finalizada y conservará el recorrido realizado.

\capturamanual{figures/seguimiento-sesion-movil.png}{Seguimiento y cierre de una sesión de muestreo.}{0.42}

\subsection{Registrar una observación}

Una observación se crea desde una sesión o mediante la opción de observación rápida.
El usuario debe escoger la campaña y la sesión correspondientes, seleccionar el tipo
de residuo y completar la información observada.

\capturamanual{figures/creacion-observacion-movil.png}{Registro de observación en campo.}{0.42}

Para elegir el tipo de residuo, la app muestra un selector con buscador, favoritos y
categorías. El usuario puede buscar por nombre o código y seleccionar el elemento que
mejor describa el residuo encontrado.

\capturamanual{figures/seleccion-residuo-movil.png}{Selección del tipo de residuo.}{0.42}

Después se indica la cantidad y, cuando proceda, detalles como tamaño, color o notas.
La app permite añadir fotografías desde la cámara o desde la galería, con un máximo
de cinco fotos por observación. Si se quieren registrar varios residuos seguidos,
puede usarse la opción \textit{Guardar y añadir otra}.

\capturamanual{figures/creacion-observacion-movil2.png}{Fotografías asociadas a una observación.}{0.42}

\subsection{Mapa móvil}

El apartado \textit{Mapa} permite consultar las campañas, sesiones, observaciones,
recorridos y áreas muestreadas desde el dispositivo. También muestra la última
posición conocida cuando la ubicación está disponible.

\capturamanual{figures/mapa-movil.png}{Mapa de trabajo en campo.}{0.42}

Este mapa es especialmente útil durante una sesión activa, porque ayuda al usuario a
mantener la orientación y comprobar por dónde se ha realizado el muestreo.

\subsection{Sincronización}

La pantalla \textit{Sincronización} muestra si existen datos pendientes de subida:
grupos, campañas, sesiones, observaciones o fotografías. Cuando el dispositivo vuelve
a tener conexión, el usuario puede pulsar \textit{Sincronizar ahora}. Si la
sincronización automática está activada, la app intentará subir los pendientes al
recuperar conexión.

\capturamanual{figures/sincronizacion-movil.png}{Sincronización de datos pendientes.}{0.42}

Durante la sincronización se muestra el progreso. Si el usuario ha configurado la
opción de sincronizar solo por Wi-Fi, la subida se bloqueará hasta que el dispositivo
esté conectado a una red Wi-Fi.

\subsection{Configuración móvil}

Desde el perfil se accede a \textit{Configuración}. En este apartado el usuario puede
cambiar el idioma de la app, activar o desactivar la sincronización automática,
indicar si quiere sincronizar solo con Wi-Fi y activar el uso de ubicación precisa.

\capturamanual{figures/configuracion-movil.png}{Preferencias locales de la aplicación móvil.}{0.42}

Estas preferencias se guardan en el dispositivo. Cambiarlas no afecta a otros
usuarios ni al portal web, salvo el idioma de la interfaz cuando se use la misma
cuenta en ese dispositivo.

\section{Trabajo sin conexión}

OmniLitter permite seguir trabajando aunque se pierda la conexión durante el trabajo
de campo. En ese caso, la aplicación móvil guarda los datos localmente y los marca
como pendientes de sincronización.

Para reducir problemas, se recomienda preparar la salida de campo con estos pasos:

\begin{enumerate}[leftmargin=1.5em]
  \item Iniciar sesión con conexión antes de salir.
  \item Seleccionar el grupo activo correcto.
  \item Abrir el apartado de campañas para cargar la información del grupo.
  \item Comprobar que las campañas y sesiones necesarias aparecen en la app.
  \item Revisar que el dispositivo tiene permisos de ubicación y cámara.
\end{enumerate}

Si se crean datos sin conexión, estos seguirán disponibles en el dispositivo. No
aparecerán en el portal web hasta que la aplicación se sincronice correctamente. Por
ese motivo, después de una jornada de campo conviene abrir \textit{Sincronización} y
confirmar que no quedan elementos pendientes.

\section{Recomendaciones de uso}

\begin{itemize}[leftmargin=1.5em]
  \item Comprobar siempre el grupo activo antes de crear campañas, sesiones u
  observaciones.
  \item Registrar las observaciones desde la sesión correcta para mantener el
  contexto del muestreo.
  \item Añadir fotografías cuando ayuden a validar o interpretar el residuo.
  \item Finalizar la sesión al terminar el recorrido.
  \item Revisar la pantalla de sincronización al volver a tener conexión.
  \item Usar los filtros del portal web antes de exportar datos, para descargar
  exactamente el conjunto que se desea analizar.
\end{itemize}

\section{Problemas frecuentes}

\begin{description}
  \item[No aparecen campañas u observaciones.] Revisar que el grupo activo es el
  correcto. En el portal, el selector está en el menú lateral; en la app móvil, en
  el perfil.
  \item[No puedo crear o editar una campaña.] La acción está reservada a
  administradores de grupo o superadministradores. Si la opción no aparece, el
  usuario solo tiene permisos de consulta.
  \item[La app no deja enviar una invitación.] Las invitaciones necesitan conexión a
  internet.
  \item[La ubicación no se guarda.] Comprobar que el dispositivo tiene el GPS activo
  y que OmniLitter tiene permiso para acceder a la ubicación.
  \item[Las fotos no se añaden.] Revisar los permisos de cámara o galería. Cada
  observación admite hasta cinco fotografías.
  \item[Los datos de campo no aparecen en el portal.] Entrar en
  \textit{Sincronización} desde la app móvil y comprobar si quedan elementos
  pendientes.
  \item[No puedo iniciar sesión sin conexión.] La cuenta debe haberse usado al menos
  una vez con conexión en ese dispositivo.
\end{description}
